% =========================================================================
% Section 7 — Conclusion
% Fragment for \input from main.tex (no preamble).
% =========================================================================
\section{결론}
\label{sec:conclusion}

본 논문은 register 를 갖춘 frozen DINOv2 ViT-B/14 인코더(86\,M 파라미터)를
14.6\,M 학습 가능 파라미터의 DPT 스타일 다중 스케일 decoder 와 결합한,
powder bed fusion 적층제조를 위한 레이어 단위 결함 세그멘테이션 모델
DefectSegAM 을 제시하였다. 모델은 레이어당 두 카메라 뷰---용융
직후(post-melting) 이미지와 도포 직후(post-spreading) 이미지---를 모두
입력받아, 두 피처 맵과 그 차분을 채널 방향으로 연결(concatenation)하는
명시적 변화 단서(change cue)를 통해 네 개의 transformer 깊이에서 융합한
뒤, $1036 \times 1036$ 해상도의 11-class 픽셀 단위 결함 맵을 디코딩한다.
라벨 효율적인 2-stage 레시피는 먼저 production DSCNN 이 의사 라벨링한
12{,}607개 ORNL 빌드 레이어로 decoder 를 사전학습한 뒤, laser powder bed
fusion 과 binder jetting 에 걸친 여덟 개 장비/소재 소스에서 추출한 단
125개의 전문가 주석 레이어로 미세조정하며, 역빈도 class-aware 샘플링,
희소 결함 클래스의 Copy-Paste 증강, 그리고 focal loss 를 통해 15.33
billion 개의 powder 픽셀에서 가장 희소한 결함 클래스의 25 픽셀까지 이르는
극단적인 클래스 불균형에 대응한다. 아울러 실험과는 별도로, 데이터셋 어휘에
혼재된 원인 기준과 외형 기준을 분리하여 레이어 이미지에 나타난 특징을 판정
기준으로 정의한 결함 8종(Under Melting 처리 확정 전 기준)과 배경 2종의
분류 체계를 향후 주석을 위해 제안하고 기존 라벨 공간 및 ASTM E3353-22
표준 용어와의 대응을 명시하였다.

ORNL Co-Registered 빌드에서 DefectSegAM 은 production 모델이 신뢰할 만한
곳에서는 이를 가깝게 따르면서도, 일관되고 결함에 민감한 방향으로
이탈한다. 결함 다수(defect-rich) 빌드 B1.2 에서 이전 내부 세대의
0.4291--0.4455 대비 0.5047 의 일치도(agreement) mIoU 를 달성하며, 이는
lack-of-fusion 및 keyhole 기공의 전조인 용융 에너지 클래스 Under
Melting(0.8643)과 Over Melting(0.4715)이 주도한 결과이다. 이 맵들의
가치는 세그멘테이션과 전혀 별개인 외부 태스크에서 입증되었다. 6{,}373개
시편에 대한 그 외 조건이 동일한 segmentation-only 인장 물성 파이프라인에서
데이터셋 제공 DSCNN 맵을 DefectSegAM 으로 교체하자, 인장강도(UTS) RMSE 가
naive 전역 평균 기준선 68.4\,MPa 보다도 나쁜 70.7\,MPa 에서 37.5\,MPa 로
감소하였으며, 항복강도(YS), 인장강도(UTS), 균일연신율(UE),
총연신율(TE) 전반에 걸쳐 해당 기준선 대비 26--45\,\% 의 일관된 감소를
보였다. 이 붕괴-복원(collapse-to-recovery) 대비는 다운스트림 물성 예측을
지배하는 것이 집계 표현의 단순성이 아니라 세그멘테이션 소스의 정보량임을
보여준다. 결정적인 세그멘테이션 수준 비교---2025 Peregrine 릴리스의
held-out 사람 ground truth 와 수집이 진행 중인 저자 라벨링 Veltz
M160/M280 및 EOS 이미지에서 충실히 재구현된 DSCNN 을 상대로 하는
비교---는 본 논문에서 완전하게 명세되었으며 그 수행은 향후 과제로 남아
있다. \TODO{\ref{sec:results}절에 스캐폴딩된 held-out, cross-machine,
ablation 평가를 완료할 것.}

여러 방향이 자연스럽게 이어진다. 첫째, 실시간 배포이다. frozen 인코더
설계는 체크포인트당 14.6\,M 학습 가능 파라미터만 저장하지만, 세그멘테이션이
생산 장비의 도포(recoating) 간격 내에 완료되도록 레이어당 추론 지연시간을
벤치마크하고 필요하다면 단축해야 한다. 둘째, 폐루프(closed-loop) 제어이다.
빌드 도중 가용한 결함 맵은 재도포(re-coating)나 파라미터 조정과 같은 교정
조치를 촉발하여 모니터를 액추에이터로 전환할 수 있다. 11개 클래스 중 어느
것이 이러한 개입을 게이트할 만큼 신뢰할 수 있는지를 정량화하는 것은
held-out 평가가 답하기 시작할 열린 질문이다. 셋째, 범위의 확장이다. 전문가
주석 코퍼스를 현재의 125개 레이어 너머로---특히 도메인 갭이 가장 큰
binder jetting 과 추가적인 합금·장비에 대해---확장하고, DefectSegAM 을
SS~316L 이외 소재의 물성 예측 파이프라인에 결합하는 것은 본 연구에서
입증된 라벨 효율성이 계속 증가하는 powder bed fusion 시스템의 보급 기반
전반에 걸쳐 누적적으로 작동하는지를 시험하게 될 것이다. 넷째, 분류 체계의
확정이다. \ref{sec:method-taxonomy}절의 외형 기준 체계로 재주석하고
재학습하면 Swelling 과 Super-Elevation 처럼 외형이 유사하여 구별하기 어려운
클래스 쌍은 하나로 합쳐지지만, 신설되는 Powder Collapse 와 ORNL 코퍼스에
25픽셀뿐인 Recoater Hopping 은 여전히 신규·희소 주석에 의존한다. 또한
기존 Over Melting 라벨은 새 정의를 만족하는지 레이어별로 재검토한 뒤
재매핑해야 하고, Under Melting 의 처리도 확정되어야 한다. 체계의 효과는
주석자 간 일치도와 재산출된 per-class 일치도로 검증해야 하며, 이는 향후
과제로 남는다.
